\documentclass[11pt,a4paper]{article}

\usepackage[utf8]{inputenc}
\usepackage[T1]{fontenc}
\usepackage{lmodern}
\usepackage{microtype}
\usepackage[margin=1in]{geometry}
\usepackage{booktabs}
\usepackage{tabularx}
\usepackage{array}
\usepackage{xcolor}
\usepackage{amssymb}
\PassOptionsToPackage{hyphens}{url}
\usepackage{hyperref}
\usepackage{graphicx}

\setlength{\emergencystretch}{3em}
\widowpenalty=10000
\clubpenalty=10000

\hypersetup{
  colorlinks=true,
  linkcolor=black,
  citecolor=blue!50!black,
  urlcolor=blue!50!black,
  pdftitle={RayaChain --- The Intelligence Layer for the Internet of Value},
  pdfauthor={Basel Magableh},
}

\newcommand{\yes}{\textcolor{green!55!black}{\textbf{\checkmark}}}
\newcommand{\no}{\textcolor{red!70!black}{\textbf{\textsf{X}}}}
\newcommand{\partialmark}{\textcolor{orange!85!black}{\textbf{\textasciitilde}}}

\title{RayaChain: The Intelligence Layer\\
       for the Internet of Value}
\author{%
  Basel Magableh\textsuperscript{1,*} \quad\quad RayaChain Research\textsuperscript{2} \\[0.8em]
  \normalsize\textsuperscript{1}School of Computer Science, Technological University Dublin, Ireland \\
  \normalsize\textsuperscript{2}Rayachain Lab, Dublin, Ireland \\[0.4em]
  \normalsize\textsuperscript{*}Correspondence: \texttt{basel.magableh@tudublin.ie}
}
\date{May 2026 \\ \textbf{v0.5 (positioning draft)}}

\begin{document}
\maketitle

\begin{abstract}
The Internet of Value moves capital, but not the context that capital
needs to be safe: counterparty reputation, route health, sentiment,
liquidity stress, and the policy a system is willing to commit to.
Bridges, oracles, and interoperability networks transport bytes and
balances across chains; none of them propagate evaluated risk. This
paper positions \emph{RayaChain} as the missing intelligence layer for
that fabric, and differentiates it on five axes ---
cross-chain coordination, semantic risk propagation, AI/agentic
reasoning, predictive infrastructure, and portable financial context ---
against LayerZero, Chainlink~CCIP, Axelar, Bittensor, and
EigenLayer/EigenCloud. We anchor the position in artefacts already on
mainnet: the Solana SPL canonical mint, six EVM OFT contracts at a
deterministic CREATE2 address, a registry-v3 commit recording the
canonical-origin flip to Solana mainnet, a continuously-running
liquidity-defence agent, and a Bittensor verification subnet with
5{,}097 successful soak rounds. We are explicit about what is not yet
shipped and list the five deliverables that close the gap between the
narrative and the visible product.
\end{abstract}

\section*{The thesis}
\label{sec:thesis}

The reframe is small in words and load-bearing in consequence: the
Internet of Value will not be safe until \emph{capital carries
evaluated risk and credit context with it across chains}. Today, value
moves freely; intelligence does not. A token bridged from Ethereum to
Solana arrives stripped of every signal that mattered on the way out
--- the route the swap took, the sentiment around the asset that
hour, the liquidity stress on the destination venue, the reputation of
the receiving wallet, the policy the sending protocol committed to.
Bridges deliver bytes and balances. They do not deliver \emph{meaning}.

RayaChain is the intelligence layer for that fabric. The unique angle
is not bridge, not oracle, not interoperability network: it is the
substrate that propagates \emph{semantic} risk and credit signals
alongside the value flows themselves. Five differentiators define the
position taken in this paper. (1)~Cross-chain coordination of risk
state, not just message transport. (2)~Semantic risk
propagation --- signals that mean something to a downstream protocol,
not raw oracle reads. (3)~AI / agentic reasoning over those signals,
constitutionally constrained so the agent cannot exceed a human-set
policy. (4)~Predictive infrastructure (forecasts of price, liquidity,
volatility, route health) rather than only post-hoc reporting.
(5)~Portable financial context --- wallet reputation and credit
posture that survives a cross-chain transfer. Each axis is necessary;
the combination is what no project listed in this paper composes
today.

\section*{Where everyone else stops}
\label{sec:comparators}

Each comparator below has shipped infrastructure we respect, and each
covers \emph{some} of the five axes. None covers all five, and most
were not designed to.

\paragraph{LayerZero.} Closest on value/message transport. LayerZero
moves arbitrary cross-chain messages, including OFT-format token
transfers, between dozens of chains. Every cross-chain RYA transfer in
this paper rides LayerZero's V2 endpoints; we depend on it operationally
and credit it as part of our substrate. What LayerZero does not do is
make the messages \emph{mean} something --- packets are opaque to the
relayer and the protocol that receives them. RayaChain adds the
semantic and risk-aware layer above the transport, so a downstream
contract sees not only ``X tokens arrived'' but ``X tokens with
sentiment $s$, route health $h$, and reputation $r$.''

\paragraph{Chainlink CCIP.} A general cross-chain communication layer
plus oracle messaging. CCIP is the strongest competitor on
\emph{infrastructure transport reliability} and on the oracle adjacency
that lets a developer attach a price read to a cross-chain message.
What CCIP is not is predictive AI risk coordination: a CCIP message
carries data but does not reason about the data, and the surrounding
oracle network is read-oriented, not policy-oriented. RayaChain treats
the oracle read as one input among five, and acts on the composed
signal under a constitutional policy.

\paragraph{Axelar.} A cross-chain interoperability network with a proof
of stake validator set that brokers connectivity and routing across a
large surface of chains. Axelar's strength is generality of
connectivity. It is not an agentic risk-intelligence layer: it does
not produce forecasts, does not reason about counterparty reputation,
and does not propagate semantic context. The relationship is the same
as with LayerZero --- Axelar can be one of the transports beneath
RayaChain, but the intelligence layer is a different problem.

\paragraph{Bittensor.} Philosophically the closest. Bittensor builds a
distributed, incentive-aligned intelligence economy: subnets compete
to produce validated outputs, and stake-weighted scoring rewards
miners that are right and consistent. RayaChain reuses that
verification primitive: our Bittensor subnet (currently shadowing on
testnet) is the validator-loss substrate for cross-chain risk scoring.
What Bittensor lacks for our purposes is a \emph{cross-chain financial
signaling layer}; subnets verify generic miner outputs but do not
themselves coordinate value-flow risk between Ethereum, Solana, and
Bitcoin-adjacent venues. RayaChain composes Bittensor's verification
with the cross-chain transport so that the validated signal lands on
the chain where the capital is.

\paragraph{EigenLayer / EigenCloud.} Trust markets, restaking, and
shared-security verification. EigenLayer makes restaked ETH economic
collateral for arbitrary Actively Validated Services; EigenCloud
extends this towards verifiable off-chain compute. This is the most
plausibly adjacent project: \emph{could become adjacent eventually},
in the CEO's framing. A future RayaChain could publish its scoring
verification as an AVS, paid for in restaked stake. We do not claim
that integration today; we note the seam.

\section*{The 5-axis differentiation matrix}
\label{sec:matrix}

Table~\ref{tab:matrix} compares the five comparators and RayaChain on
the five axes from \S\,\nameref{sec:thesis}. Cells are \yes\ (shipped and
load-bearing for the project), \partialmark\ (partial --- adjacent or
in-progress), or \no\ (out of scope for the project). Each cell is a
one-line justification, not an opinion. RayaChain's row is calibrated
to what is on mainnet today (\S\,\nameref{sec:mainnet}); axes that depend on
the operational fabric are marked partial accordingly
(\S\,\nameref{sec:gap}).

\renewcommand{\arraystretch}{1.25}
\begin{table}[ht]
\centering
\small
\newcolumntype{Y}{>{\raggedright\arraybackslash}X}
\begin{tabularx}{\textwidth}{@{}l Y Y Y Y Y@{}}
\toprule
 & \textbf{Cross-chain coord.} & \textbf{Semantic risk prop.}
 & \textbf{AI / agentic reasoning} & \textbf{Predictive infra.}
 & \textbf{Portable financial context} \\
\midrule
\textbf{LayerZero}
 & \yes\ V2 endpoints move messages and OFT tokens across many chains
 & \no\ payloads are opaque to the protocol; no risk semantics
 & \no\ relayers do not reason
 & \no\ no forecasts emitted
 & \no\ token state ports; reputation does not \\
\textbf{Chainlink CCIP}
 & \yes\ general cross-chain messaging plus oracle adjacency
 & \partialmark\ price feeds attach as data, not as risk signals
 & \no\ no agentic reasoning layer
 & \partialmark\ feeds are forecasts of \emph{prices}, not route health or contagion
 & \no\ no portable reputation primitive \\
\textbf{Axelar}
 & \yes\ PoS validator set brokers connectivity and routing
 & \no\ payloads carry no risk semantics
 & \no\ no reasoning layer
 & \no\ no forecasts emitted
 & \no\ no portable reputation primitive \\
\textbf{Bittensor}
 & \no\ subnets do not coordinate cross-chain financial state
 & \partialmark\ subnet outputs can encode semantics but not financial signaling
 & \yes\ stake-weighted incentive layer over miner reasoning
 & \partialmark\ subnets exist for forecasting, but not as a cross-chain primitive
 & \no\ no cross-chain wallet reputation primitive \\
\textbf{EigenLayer / EigenCloud}
 & \partialmark\ AVSs can be cross-chain but EigenLayer itself is L1-anchored
 & \no\ shared-security primitive, not a semantics primitive
 & \no\ no agentic reasoning layer
 & \no\ no forecasts emitted
 & \partialmark\ restaked-stake reputation, not portable wallet credit \\
\midrule
\textbf{RayaChain}
 & \yes\ canonical Solana mainnet mint plus six EVM OFT peers at one CREATE2 address
 & \partialmark\ semantic signals encoded; live cross-chain propagation pending (\S\,\nameref{sec:gap})
 & \yes\ constitutional agent (RAYA defence bot) firing under written policy since 2026-05-06
 & \partialmark\ five-engine architecture shipped; live downstream consumers pending
 & \partialmark\ wallet-reputation primitive specified; portability across chains is the next gap \\
\bottomrule
\end{tabularx}
\caption{Five-axis differentiation matrix. Justifications are
one-liners; deeper architectural claims live in the canary paper and
the long technical report (links in \S\,\nameref{sec:read-more}).
RayaChain's partial cells are the same five deliverables enumerated in
\S\,\nameref{sec:gap}.}
\label{tab:matrix}
\end{table}

The table reads cleanly along two diagonals. Horizontally, no
comparator covers all five axes. Vertically, the axes that no
non-RayaChain project marks \yes\ are \emph{semantic risk propagation}
and \emph{portable financial context} --- the two axes that follow
from the thesis that capital should carry context. Those are the axes
we own.

\section*{What is shipped on mainnet}
\label{sec:mainnet}

The CEO directive that motivates this paper is also a constraint:
\emph{all features are implemented on mainnet}. Every claim below
traces to a public artefact a reader can verify themselves.

\begin{itemize}
  \item \textbf{Solana mainnet RYA SPL mint}, live since 2026-05-04:%
    \\[0.15em]
    \mbox{\texttt{EebqvpbYoJYQLzeUfjVoHWopogSE6ECRX5HJmBr6susV}}.
    Defended by the OmniRisk hot-wallet bot since 2026-05-06;
    explorer-verifiable on Solscan.
  \item \textbf{EVM OFT contracts} at the deterministic CREATE2
    address%
    \\[0.15em]
    \mbox{\texttt{0x1a281E61ed9f16D325D93d50a50D008b3d141676}}
    on Ethereum, Base, Polygon, Optimism, and BSC mainnet
    (\texttt{name()=RayaChain}, \texttt{symbol()=RYA} verified on
    each chain's explorer); the same CREATE2 address is inferred for
    Arbitrum mainnet pending re-verification --- 5 of 6 chains
    explorer-confirmed today.
  \item \textbf{Registry v3} (committed 2026-05-08, OmniRisk
    issue OMN-181) records the canonical-origin flip from
    \texttt{solana-devnet} to \texttt{solana-mainnet} and lists eight
    mainnet entries, making the on-disk source of truth match the
    on-chain reality.
  \item \textbf{RAYA defence bot}: a \emph{buy-only}
    liquidity-defence agent operating under written policy since
    2026-05-06. The bot is the agentic-engine demonstration; the
    cross-chain canary layer is the substrate it sits on. The
    canary paper documents the 2026-05-08 to 2026-05-09 defence
    cycle as a case study, not as a market-manipulation product.
  \item \textbf{Bittensor verification subnet} (shadow deployment on
    testnet metagraph): \textbf{5{,}097} successful rounds with
    \textbf{0} authentication failures. This validates layer-1
    orchestration; layer-2 economic verification under real metagraph
    stakes is deferred per the canary paper's honesty split, and is
    the fifth deliverable in \S\,\nameref{sec:gap}.
\end{itemize}

These five artefacts are the empirical floor of the position. Nothing
above the floor depends on a fact we cannot point at.

\section*{Narrative vs visible product}
\label{sec:gap}

To use the CEO's phrase verbatim: \emph{the narrative is ahead of the
visible product}. We say this not to retreat from the position but to
mark it precisely. The contracts have shipped; the operational fabric
that turns shipped contracts into visible cross-chain risk
intelligence has not yet closed. This is the deliberate engineering
order of operations, not an oversight.

Five deliverables close the gap. None is speculative; each has an
issue and an owner.

\begin{enumerate}
  \item \textbf{Live cross-chain risk-signal propagation}, not just
    SPL transfers. A risk message published on the canary substrate
    must arrive at a downstream chain and be readable as
    \emph{semantics}, not as opaque bytes.
  \item \textbf{Predictive alerts consumed by downstream protocols},
    not just rendered on the OmniRisk dashboard. The forecast must
    cross a protocol boundary and change a third party's behaviour.
  \item \textbf{Portable wallet-reputation across cross-chain
    transfers.} A wallet's credit posture on Ethereum must survive
    its appearance on Solana --- the load-bearing primitive for
    \emph{capital carrying context}.
  \item \textbf{Third-party audit} of the EVM OFT contracts and the
    Bittensor scoring code. Internal verification is not enough for
    the trust boundary the system claims.
  \item \textbf{Subtensor mainnet deployment} of the verification
    subnet. The 5{,}097-round soak demonstrates layer-1 correctness;
    economic verification under real metagraph stakes is what closes
    the loop.
\end{enumerate}

We are not asking the reader to take any of the five on faith. We are
asking the reader to read this paper as the position taken at the
contract-shipped, fabric-pending boundary --- because that is exactly
where the project is.

\section*{What this paper is not claiming}
\label{sec:not-claiming}

Three disclaimers, in plain English.

\begin{enumerate}
  \item \textbf{Not Bittensor-style economic verification under real
    metagraph stakes today.} The verification subnet has soaked on
    testnet (5{,}097 rounds, 0 auth failures); a Subtensor mainnet
    deployment under real stake is deferred and explicit, per the
    canary paper's layer-1/layer-2 honesty split.
  \item \textbf{Not a finished operational fabric.} The contracts
    are deployed. Peer-wiring across all six EVM chains, the executor
    sidecar, and the third-party audit are work in progress and are
    explicitly listed in \S\,\nameref{sec:gap}. Anything above ``contracts
    deployed and registry coherent'' is being closed, not claimed
    closed.
  \item \textbf{Not a generic agentic-engine product.} The RAYA
    defence bot is a single-token, buy-only liquidity defence run
    under a written policy. It is the agentic-engine demonstration on
    mainnet; it is not a product offered to other protocols today,
    and the bot's policy is not a substitute for protocol-level risk
    intelligence. The product surface is the intelligence layer; the
    bot is one of its first concrete users.
\end{enumerate}

\section*{Read more}
\label{sec:read-more}

This positioning paper deliberately leaves the technical depth to its
companion documents.

\begingroup
\sloppy
\hbadness=10000
\begin{itemize}
  \item \textbf{Canary paper} --- \emph{The RayaChain Omnichain
    Canary: Validator-Loss Verification for Cross-Chain Risk
    Signaling}. The architectural source of the layer-1/layer-2
    honesty split and the validator-loss decomposition. URL:
    \url{https://green.omnirisk.io/media/papers/2026/rayachain-omnichain-canary/manuscript.pdf}.
  \item \textbf{Long technical report} --- \emph{OmniRisk $\times$
    RayaChain $\times$ Bittensor: An End-to-End Risk-Intelligence
    Stack for the Internet of Value}. The full mainnet implementation
    record (\S17), and the migration plan to multisig ownership.
    URL:
    \url{https://green.omnirisk.io/media/papers/2026/omnirisk-rayachain-bittensor/manuscript.pdf}.
  \item \textbf{Focused architecture paper} --- \emph{Agentic,
    Context-Aware Risk Intelligence in the Internet of Value}. The
    five-engine architecture and the constitutional-agent rule.
    URL:
    \url{https://green.omnirisk.io/media/papers/2026/agentic-iov-risk-focused/manuscript.pdf}.
  \item \textbf{RYA tokenomics} --- on Paragraph at
    \url{https://paragraph.com/@omnirisk}. This paper does not
    restate tokenomics.
  \item \textbf{On-chain artefacts.} Solana mint on Solscan:%
    \\[0.15em]
    \href{https://solscan.io/token/EebqvpbYoJYQLzeUfjVoHWopogSE6ECRX5HJmBr6susV}%
         {\texttt{solscan.io/token/Eebqvpb\dots susV}}.
    EVM OFT (Ethereum, Base, Polygon, Optimism, BSC, and inferred
    Arbitrum) at CREATE2 address
    \mbox{\texttt{0x1a281E61ed9f16D325D93d50a50D008b3d141676}};
    Etherscan / Basescan / Polygonscan / Optimistic Etherscan /
    BscScan listings carry the deployment.
\end{itemize}
\endgroup

\noindent\textit{Status note:} this paper is the v0.5 first populated
draft per OMN-183 (2026-05-09), pending CEO sign-off on venue. No
external publication has been made.

\end{document}
